\chapter{Wstęp}
\section{Wprowadzenie do tematyki}
Systemy monitoringu wizyjnego (CCTV) stanowią obecnie kluczowy element utrzymywania bezpieczeństwa w miastach. Większość systemów monitoringowych polega na zapisie pełnego strumienia wideo w celu późniejszej ewentualnej analizy, w przypadku wystąpienia jakiegoś incydentu. Z biegiem lat i wraz ze wzrostem liczby instalowanych kamer, takie podejście staje się coraz trudniejsze, ze względu na czasochłonność ewentualnej analizy wideo z wielu perspektyw.

Przełomem w tej dziedzinie stał się rozwój algorytmów sztucznej inteligencji, a w szczególności ich wykorzystania w widzeniu komputerowym (ang. \textit{Computer Vision}). Współczesne systemy analityki wideo potrafią w czasie rzeczywistym klasyfikować obiekty, śledzić ich ruch i rozpoznawać anomalie. W tę właśnie dziedzinę wpisuje się tematyka niniejszego projektu, którego celem jest przeniesienie problemu detekcji niebezpiecznych lub nietypowych zachowań z człowieka na zautomatyzowany system działający bezpośrednio na strumieniach wideo w czasie rzeczywistym.


\section{Analiza problematyki}
Pomimo powszechności systemów kamer, klasyczne podejście do monitoringu napotyka na szereg istotnych ograniczeń natury ludzkiej, prawnej oraz technologicznej. Projekt rozwiązuje te problemy z następujących obszarach:

\begin{itemize}
\item \textbf{Ograniczenia ludzkiej percepcji} --- tradycyjne podejście do monitoringu polega na ciągłej obserwacji kilkunastu lub kilkudziesięciu ekranów przez operatora. Badania z zakresu ergonomii i bezpieczeństwa dowodzą, że po upływie zaledwie 20 minut ciągłej obserwacji zdolność do wychwycenia istotnego zdarzenia drastycznie spada~\cite{human_factor_cctv}. Z kolei generowanie powiadomień na podstawie sztywnego wykrywania ruchu może prowadzić do ogromnej liczby fałszywych alarmów. Z tego względu zauważalna jest potrzeba implementacji logiki opartej na elastycznych szablonach, które alarmują użytkownika dopiero w momencie zajścia specyficznej kombinacji zdarzeń (np.\ pojawienie się w kadrze niebezpiecznego narzędzia i jednoczesny upadek człowieka).
\item \textbf{Ochrona prywatności} --- przechowywanie surowych nagrań wideo z przestrzeni publicznych wiąże się z dużym ryzykiem naruszenia prywatności przechodniów. Ogólne rozporządzenie o ochronie danych (RODO) narzuca zasadę minimalizacji gromadzonych danych~\cite{gdpr_video_surveillance}. Rejestrowanie wizerunków twarzy przechodniów lub numerów rejestracyjnych pojazdów stanowi problem prawny. Stąd wynika konieczność integracji modułów anonimizujących dane jeszcze przed ich zapisem w bazie danych.
\item \textbf{Zarządzanie długimi nagraniami wideo} --- ciągły zapis strumieni wideo, nawet w niskiej rozdzielczości, generuje gigantyczne ilości danych. Najprostsze systemy zapisują materiały w postaci dużych, jednolitych plików, co rodzi problemy na etapie ich późniejszego odczytu ze względu na czasochłonność buforowania i przeszukiwania ogromnych bloków danych. W związku z tym optymalizacja procesu zapisu, na przykład poprzez rejestrację jedynie fragmentów zawierających wykrycia oraz ich fragmentację na krótkie segmenty czasowe z wykorzystaniem standardów takich jak \texttt{MPEG-DASH}~\cite{dash_standard}, stanowi jedno z kluczowych zagadnień projektowych.
\end{itemize}